\section{Processing transparency}\label{sec:5}

Section 4 measured rated transparency at the level of whole compounds; the rating is a steady-state property of the compound, given by speakers' explicit judgements. Section 5 turns to a time-course property: how does semantic transparency modulate the processing of morphologically complex words during recognition? Heyer and Kornishova's \autocite*{heyer-kornishova-2018-priming-eventually} masked-priming study is the empirical anchor. Their finding, summarised in their title's three-dot ellipsis: semantic transparency affects morphological priming, but only at long stimulus-onset asynchronies. §\ref{sec:5-1} lays out the design and findings. §\ref{sec:5-2} unpacks the serial-processing interpretation. §\ref{sec:5-3} places it against a rival camp. §\ref{sec:5-4} returns to the cluster question. §\ref{sec:5-5} closes.

\subsection{The Heyer and Kornishova study}\label{sec:5-1}

Heyer and Kornishova \autocite*{heyer-kornishova-2018-priming-eventually} ran two parallel masked-priming experiments, one with English -\mention{ness} nominalisations and one with Russian -\mention{ost'} nominalisations. The two suffixes are deadjectival and produce abstract nouns denoting ``state of being $X$'' (e.g., \mention{dark + ness = darkness}; \mention{bledn-} + \mention{-ost'} = \mention{blednost'} ``paleness''). Targets in both experiments were the adjective bases (e.g., \mention{pale, dark}); primes were the corresponding nominalisation, a simplex unrelated control, or the adjective as identity prime. Each experiment crossed two stimulus-onset asynchronies (SOAs), short and long, with transparency as a continuous variable rated by separate groups of speakers (Appendix C).

Transparency was rated on a 1--7 scale: \mention{business} (rated 3.17) at the opaque end is no longer interpretable as ``state of being busy''; \mention{falseness} (rated 5.75) is fully derivable from \mention{false}; items spanned the range continuously.

The headline result is SOA-conditioned. At short SOA in both languages, priming is constant across the transparency scale: morphologically complex primes are decomposed into stems regardless of whether the stem still contributes to the meaning of the whole. At long SOA, priming becomes transparency-modulated: more transparent prime-target pairs (e.g., \mention{paleness-pale}) yield larger priming effects than less transparent pairs (e.g., \mention{business-busy}). The three-way interaction (Prime Type × Transparency × SOA) is significant in English and marginally so in Russian, with the same direction in both languages \autocite[1117--1119]{heyer-kornishova-2018-priming-eventually}. Statistical detail is in Appendix C.3.

\subsection{The serial-processing interpretation}\label{sec:5-2}

Heyer and Kornishova read the data as evidence for a serial decomposition process. Morphologically complex words are first decomposed into stem-plus-affix on the basis of orthographic structure (\mention{morpho-orthographic decomposition}), with semantic information entering only at a later stage. The first stage is semantically blind: it parses \mention{paleness} into \mention{pale + ness} regardless of whether the parse yields a coherent semantic interpretation, and would parse \mention{corner} into \mention{corn + er} even though the parse fails semantically \autocite[1112]{heyer-kornishova-2018-priming-eventually}.

This is consistent with Bell and Schäfer's \autocite*{bell-schafer-2016-modelling-transparency} §\ref{sec:4} hypothesis that ``perceived transparency may itself be a reflex of ease of processing.'' Within the narrow Heyer-and-Kornishova-style derived-word masked-priming design, the pattern is that fast and shallow processing windows don't show transparency modulation, while longer processing windows do. The ``fast-and-shallow versus slow-and-deep'' framing applies to that paradigm specifically; whether it extends to other tasks (whole-word reading, lexical decision without priming, comprehension) is a separate empirical question. The two studies are compatible with a shared form-meaning accessibility hypothesis, though they involve different items (NN compounds versus derivational nominalisations) and different methods (ratings/modelling versus masked priming) and so don't test the same processing chain.

\subsection{The rival parallel-processing position}\label{sec:5-3}

The serial reading is contested. Feldman and colleagues' parallel-immediate position argues that morpho-semantic information influences recognition from the earliest stages. Feldman et al.'s \autocite*{feldman-etal-2015-must-analysis} time-course analysis presses the diagnostic question (\mention{must analysis of meaning follow analysis of form?}) with a same-target design --- each target is compared after semantically similar and semantically dissimilar primes within subjects --- and SOA modelled as continuous from 34 to 100~ms rather than as binned categorical levels. They report that semantic similarity affects recognition latencies even at SOAs of 34~ms and 48~ms, shorter than Heyer and Kornishova's \mention{short} windows of 39~ms (English) and 33~ms (Russian). On the parallel-immediate reading, the form-then-meaning ordering doesn't survive that finding. The two designs aren't head-to-head: Heyer and Kornishova's short-SOA result is the absence of transparency modulation across an item set rather than the absence of semantic effects on a same-target comparison, and the disagreement is whether semantic information is constitutively present at masked-priming windows or emerges only after morpho-orthographic decomposition. The literature is mixed enough (the meta-table on p.~1121 of Heyer and Kornishova surveys 35 prior masked-priming studies) that the present paper takes the serial reading as Heyer and Kornishova's reading, not as a settled disciplinary verdict. The cross-modal probe in Appendix~\ref{sec:D} would generate within-items evidence relevant to the disagreement.

\subsection{What §\ref{sec:5} contributes to the cluster}\label{sec:5-4}

§\ref{sec:5} contributes the time-course dimension to the form-meaning accessibility family. §\ref{sec:3} framed compositional transparency as a gradient; §\ref{sec:4} made it quantitative through corpus predictors and ratings; §\ref{sec:5} makes it temporally articulated. The family's evidential support is now: a theoretical scale, a steady-state quantitative measure, and a time-course measure. The three windows are operationally distinct but conceptually related: they're different probes for the same underlying form-meaning accessibility, taken at different stages of speakers' interaction with the relevant items. Falsification condition for the within-domain priming profile: if rated semantic transparency failed to modulate masked-priming magnitude at long SOAs in an extension or replication of Heyer and Kornishova's design, the profile would fail.

The relationship between §\ref{sec:4}'s ratings and §\ref{sec:5}'s priming deserves a careful word. Bell and Schäfer's compounds and Heyer and Kornishova's nominalisations are different items in different morphological domains; their results don't compose into a single transitive empirical chain (corpus predictors $\rightarrow$ ratings $\rightarrow$ priming on the same items). What they jointly support is the cluster hypothesis: rated transparency and priming-graded transparency both behave as the form-meaning accessibility hypothesis predicts, in a way that distinguishes them from idiosyncratic post-hoc descriptive labels. They're compatible probes for the same construct, not direct measurements of one another. Whether the cluster hypothesis survives the within-items test --- corpus predictors, ratings, priming magnitude, and interpretation accuracy aligning on shared items as joint reflexes of underlying representational structure --- is the empirical question §\ref{sec:D} sketches as a probe design. Until that probe is run, the family is three within-domain profiles plus a research hypothesis, not a delivered cross-domain cluster.

\subsection{Cluster status and bridge to §\ref{sec:6}}\label{sec:5-5}

§§\ref{sec:3}--\ref{sec:5} give three within-domain projectibility profiles, each stabilised by a different source: a construction-level profile (§\ref{sec:3}, stabilised by constructional inheritance) predicts compositional interpretability via inheritance and network position; a compound-level profile (§\ref{sec:4}, stabilised by usage-based entrenchment) predicts rated transparency from corpus statistics; a derivationally-complex-word-level profile (§\ref{sec:5}, stabilised by processing architecture) predicts long-SOA priming magnitude from rated transparency. Whether they jointly license a cross-domain projectibility cluster on shared items, as opposed to remaining three within-domain profiles loosely related by analytical theme, is the question §\ref{sec:7} takes up. Acquisition trajectory is a plausible further extension, consistent with the network claim of §\ref{sec:3-3}, but the present sections don't directly ground it. Khalidi's \autocite*{khalidi-2018-causal-networks} causal-network refinement applies at the cross-domain level only if shared-item evidence supports a causal-hierarchical reading, which the present paper doesn't supply. §\ref{sec:5}'s contribution is to add a time-course profile to the within-domain evidence, not to establish cross-domain projection by itself.

§\ref{sec:6} turns from form-meaning accessibility back to the comparative question raised in §\ref{sec:1}: how does the umbrella schema constrain crosslinguistic generalisation about transparency? Haspelmath's \autocite*{haspelmath-2010-comparative-concepts} descriptive-versus-comparative-concept distinction is the right discipline for that question, and §\ref{sec:6} applies it to the morphosyntactic and form-meaning cases the paper has worked through.
